\section{Testing general higher K-theory bounds against the actual window}
\label{sec:k-theory-growth-audit}

The all-prime quantity in AC-K10 suggests using general torsion
bounds.  This section records the numerical test of that approach,
including its quantifiers.  It does not infer the size of an actual
K-group from the size of an available upper bound.

\begin{proposition}[AC-K24: the required scale of a uniform K-theory bound]
\label{AC-K24}
Use the growing fields $F_x$ of (FG.3), and put
$F_x^{\mathrm c}=F_x(\sqrt[3]{269})$ and
$t_x=\lfloor x^{401}\rfloor$.  Fix $0<\varepsilon<1/52$.
The following would suffice for the all-prime size part of
the arithmetic budget: constants $C,A>0$, independent of
$x$, the selected source primes, and $e$, such that
\[
 \log\#K_{2e}(\mathcal O_{F_x^{\mathrm c}})
 \leq C\bigl(1+e+[F_x^{\mathrm c}:\mathbb Q]
                    +\log D_{F_x^{\mathrm c}}\bigr)^A
       \quad(1\leq e<t_x).                         \tag{KG.1}
\]
Under this additional hypothesis,
\[
 \sum_{1\leq e<t_x}\sum_\ell\log\ell\,
  \operatorname{length}_{\mathbb Z_\ell}
        S_\ell(F_x^{\mathrm c},e)
       =o\bigl(\exp(x^\varepsilon/2)\bigr).         \tag{KG.2}
\]
The finite gluing errors already bounded in AC-K19 can be
added without changing (KG.2).  The actual theta/Selmer
identification, common source window, and positive budget slack
remain separate requirements.

The explicit upper-bound expression in Bayer-Fluckiger--Emery--Houriet,
Theorem~1.1, cannot certify (KG.2) by direct substitution,
even with fixed field degree and discriminant.  The newer
totally real/CM estimate of de la Torre Dur\'an has the same
numerical limitation at this weight scale, and its field
hypothesis fails for $F_x^{\mathrm c}$.
\end{proposition}
\begin{proof}
By (FG.11),
\[
 [F_x:\mathbb Q]=O_\varepsilon(x^{2+8\varepsilon}),\qquad
 \log D_{F_x}
 =O_\varepsilon\bigl(x^{2+41\varepsilon/5}(\log x)^2\bigr).
                                                               \tag{KG.3}
\]
Let $E=\mathbb Q(\omega,\sqrt[3]{269})$, a fixed field.
Since $F_x^{\mathrm c}=F_xE$, the tensor-order argument
(FG.10) gives
\[
 [F_x^{\mathrm c}:\mathbb Q]\leq6[F_x:\mathbb Q],\qquad
 \log D_{F_x^{\mathrm c}}
 \leq6\log D_{F_x}+[F_x:\mathbb Q]\log D_E.          \tag{KG.4}
\]
The sharper degree factor three is unnecessary here.
For the allowed $\varepsilon$, the right sides of (KG.3)--(KG.4)
grow more slowly than $x^4$.  Since $e<t_x\leq x^{401}$,
(KG.1) is at most $C_\varepsilon x^{401A}$ after enlarging
the constant for bounded $x$.  Summing the K10.3 inequality
over these weights gives an upper bound
$C_\varepsilon x^{401(A+1)}$ for the left side of (KG.2).
For every fixed $B>0$,
\[
       B\log x-x^\varepsilon/2\longrightarrow-\infty.
\]
This proves (KG.2).  The same calculation applies to the
polynomial error (FG.4).  It controls the arithmetic sizes
only; no step constructs a missing source comparison.

For the direct-bound test, write $m$ for the K-degree and
$d,D$ for field degree and absolute discriminant.  The
positive expression on the right of the 2015 theorem is
\[
 U(m,d,D)=(2m+1)^{71m^4d^3}
               d^{293m^5d^5}D^{528m^5d^4}.          \tag{KG.5}
\]
Take $m=2(t_x-1)$, for sufficiently large $x$.
For every $d,D\geq2$,
\[
 \log U(m,d,D)\geq528m^5d^4\log D\geq c x^{2005}.
\]
Consequently $U(m,d,D)/\exp(x^\varepsilon/2)\to\infty$.
Even this single term in a sum of the displayed upper bounds
is larger than the allowed budget.  This conclusion concerns
$U$; it is neither a lower bound for a K-group nor evidence
that (KG.1) fails.

For comparison, the 2026 bound has an envelope of the form
\[
 U_C'(m,d,D)=
 \exp\!\bigl(CN^2d(\log N+\log d)+CN^2\log D\bigr),
          \qquad N=2m+1,                            \tag{KG.6}
\]
with a fixed positive constant $C$.  At the same $m$,
even fixed $d\geq2$ gives
$\log U_C'\geq c_C x^{802}\log x$, so this displayed
envelope also does not fit the budget.  Independently of
this comparison, $F_x^{\mathrm c}$ contains $E$.
AM-H03 proves that $E$ embeds in no CM field, so
$F_x^{\mathrm c}$ cannot be CM.  It cannot be totally real,
since it contains $\omega$.  Thus it is outside the stated
field class of that theorem.
\end{proof}

\paragraph{Exact external locators and cardinality conventions.}
Bayer-Fluckiger--Emery--Houriet, \emph{Hermitian Lattices and Bounds
in K-Theory of Algebraic Integers}, Documenta Mathematica,
Merkurjev volume (2015), Theorem~1.1, p.~72, gives (KG.5)
for $\log\operatorname{card}_bK_m(\mathcal O_F)_{\rm tors}$,
$b=\max(d+1,2m+2)$.  Its introduction, p.~71, defines
$\operatorname{card}_b$ using the quotient by the subgroup
generated by elements of order at most $b$; its final
derivation is \S5.2, p.~82.
\url{https://doi.org/10.4171/DMS/7/4}.
These are inspected source statements and the displayed
final derivation, not a new proof of all their upstream inputs.

De la Torre Dur\'an, \emph{Perfect Forms and Torsion Bounds in
Algebraic K-Theory over Totally Real and CM Fields},
arXiv:2611.06300v1, \S4.2, Theorem~4.5 and its complete printed
proof, gives (KG.6) in big-O notation for the logarithm of its
auxiliary cardinality.  That paper defines this cardinality differently:
it keeps only primary parts for primes greater than $b$,
where $b=dN+1$ in the general case of \S4.2.
\url{https://arxiv.org/html/2611.06300v1#S4.SS2}.
Neither auxiliary cardinality is silently substituted for
the full order in (KG.1).

Soul\'e, \emph{A Bound for the Torsion in the K-Theory of
Algebraic Integers}, Documenta Mathematica, Kato volume
(2003), \S5.1, p.~785, formulates a fixed-degree conjecture
of shape $\#K_m(\mathcal O_F)_{\rm tors}
\leq\alpha\exp(\beta m^\gamma\log D_F)$.
The constants are degree-dependent; that conjecture as
stated does not supply (KG.1) as the degree grows.
The higher-regulator alternative is discussed in \S5.4,
pp.~786--787, equation~(26).
\url{https://doi.org/10.4171/DMS/3/21}.
No conjecture is invoked as a theorem here.

\paragraph{Geometric meaning and status.}
This is a size test for the localization and duality groups on
$\operatorname{Spec}\mathcal O_{F_x^{\mathrm c}}[1/\ell]$
identified in AC-K10.  The finite coefficient descent is the
previously audited SGA~1, V~\S\S4 and 7, and XI~\S6;
the K-theory coefficient formalism uses SGA~6, V~2.2(b).
Neither provides the uniform numerical bound (KG.1).
No Grothendieck archive scan is newly attributed a growth
estimate.  This audit records a tested approach and a
sufficient quantitative substitute; it does not close the
main infinite-coefficient depth gate or H04.4.
